\section{Runtime Monitoring}

\subsection{Resiliency Real-Time and Certification}

\begin{frame}{Runtime Monitoring - Runtime Assurance (RTA)}
    \begin{block}{Runtime Assurance (RTA)}
        \begin{itemize}
            \item \textbf{Concept:} Instead of trying to certify a highly complex or non-deterministic AI module, we wrap it in a monitor that switches to an assured recovery function upon predicting a failure.
            \item \textbf{Modularity:} The system remains certified even if the AI module is modified or replaced, because safety is guaranteed by the assured recovery module.
            \item \textbf{Versatility:} Can be applied at any level of the system architecture (e.g., controllers, actuators).
        \end{itemize}
    \end{block}
\end{frame}

\begin{frame}{Runtime Monitoring - RTA Architecture}
    \begin{columns}
        \begin{column}{.48\textwidth}
            \begin{block}{ASTM F3269-17 Standard}
                Formalized architecture for unmanned aircraft systems:
                \begin{itemize}
                    \item \textbf{Input Manager:} Validates unassured data.
                    \item \textbf{Complex Function:} Untrusted AI module performing the main task.
                    \item \textbf{Recovery Function:} Assured module that maintains a safe state.
                    \item \textbf{Safety Monitor:} Assured module that detects failures.
                    \item \textbf{Switch:} Assured routing.
                \end{itemize}
            \end{block}
        \end{column}
        \begin{column}{.52\textwidth}
            QUI IMMAGINE RTA-architecture
        \end{column}
    \end{columns}
\end{frame}

\begin{frame}{Runtime Monitoring - Hierarchical Safety Monitor}
    \begin{columns}
        \begin{column}{.52\textwidth}
            \begin{block}{Hierarchical Checking}
                The monitor can analyze the AI module hierarchically to improve safety:
                \begin{itemize}
                    \item \textbf{Low-Level Monitor:} Checks for evident, immediate errors (e.g., out-of-distribution inputs, syntax, range checks).
                    \item \textbf{High-Level Monitor:} Checks for semantic, functional, and physical constraints.
                \end{itemize}
            \end{block}
        \end{column}
        \begin{column}{.48\textwidth}
            QUI IMMAGINE hierarchical-checker
        \end{column}
    \end{columns}
\end{frame}

\begin{frame}{Runtime Monitoring - Coverage and Voter Architecture}
    \begin{columns}
        \begin{column}{.52\textwidth}
            \begin{block}{Coverage Scope}
                \begin{itemize}
                    \item \textbf{Operating Scope:} The domain of the complex function.
                    \item \textbf{Coverage Gap:} Any state where a failure cannot be detected or recovered (QUI IMMAGINE coverage-gap).
                \end{itemize}
            \end{block}
            \begin{block}{Multi-Monitor Voter}
                \begin{itemize}
                    \item Multiple monitors, each with a corresponding recovery function.
                    \item A voter chooses which monitor's recovery function to execute.
                \end{itemize}
            \end{block}
        \end{column}
        \begin{column}{.48\textwidth}
            QUI IMMAGINE monitors-voter
        \end{column}
    \end{columns}
\end{frame}
